% Journal-only tagged-cover theorem, included by the journal compression section.
% The raw module inventory is denoted U_L^{raw} in the current preprint, so the
% tagged universe below carries an explicit superscript to prevent a collision.

\subsection{An exact tagged covering representation under identity closure}
\label{subsec:aor-tagged-cover}

Fix a finite belief set \(B\), a finite module carrier \(M\), a finite
strategy catalog \(S\), and a source library \(L\subseteq S\).  Every
admissible library contains the distinguished inactive strategy \(s_0\).
Profiles are exact rational functions \(j_s:B\to\Q\),
\(j_{s_0}(b)=0\) for every \(b\), and \(\mods(s_0)=\varnothing\).  Each
\(\mods(s)\) is a subset of \(M\).  For an admissible library \(A\), let
\[
 F_A(b):=\max_{s\in A}j_s(b),\qquad
 U_A^{\rm raw}:=\bigcup_{s\in A}\mods(s),\qquad
 C_A:=\cl(U_A^{\rm raw}).
\]
The maximum is attained because \(A\) is finite and nonempty.  The inactive
strategy also implies \(F_A(b)\geq 0\).

In this subsection only, closure is the identity on \(2^M\).  Thus
\(C_A=U_A^{\rm raw}\) for every \(A\).  Define the source-relative tagged
requirement universe
\begin{equation}
 \mathcal U_L^{\rm tag}
 :=\{\mathsf{frontier}(b):b\in B\}
   \mathbin{\dot\cup}
   \{\mathsf{module}(m):m\in C_L\}.
 \label{eq:tagged-requirement-universe}
\end{equation}
The constructors are disjoint: a frontier tag and a module tag are different
even if their underlying identifiers have the same printed label.  For each
source strategy \(s\in L\), define
\begin{equation}
 \mathcal R_s
 :=\{\mathsf{frontier}(b):j_s(b)=F_L(b)\}
   \cup
   \{\mathsf{module}(m):m\in\mods(s)\}.
 \label{eq:strategy-tagged-coverage}
\end{equation}
Equality in the first set is exact.  Consequently every tied source
maximizer covers the frontier row; no tie-breaking rule or distinguished
maximizer enters the construction.

\begin{theorem}[Tagged-cover equivalence under identity closure]
\label{thm:aor-tagged-cover-equivalence}
Under the assumptions above, let \(L'\) be an admissible source sublibrary:
\(s_0\in L'\subseteq L\).  Then
\begin{equation}
 F_{L'}=F_L\ \text{and}\ C_{L'}=C_L
 \quad\Longleftrightarrow\quad
 \mathcal U_L^{\rm tag}\subseteq
       \bigcup_{s\in L'}\mathcal R_s.
 \label{eq:tagged-cover-equivalence}
\end{equation}
All belief rows are retained in \eqref{eq:tagged-cover-equivalence}.  If
\(F_L(b)=0\), the mandatory inactive strategy attains that frontier and
automatically covers \(\mathsf{frontier}(b)\).
\end{theorem}

\begin{proof}
Suppose first that \(F_{L'}=F_L\) and \(C_{L'}=C_L\).  Fix \(b\in B\).
The finite maximum defining \(F_{L'}(b)\) has an attainer
\(s\in L'\).  Frontier equality gives
\(j_s(b)=F_{L'}(b)=F_L(b)\), so
\(\mathsf{frontier}(b)\in\mathcal R_s\).  This argument is existential and
therefore permits any number of tied attainers.  Next fix \(m\in C_L\).
Identity closure and closure equality give
\[
 m\in C_{L'}=U_{L'}^{\rm raw}
   =\bigcup_{s\in L'}\mods(s).
\]
Hence some \(s\in L'\) carries \(m\), so
\(\mathsf{module}(m)\in\mathcal R_s\).  Every tagged requirement is covered.

Conversely, suppose every tagged requirement is covered.  Because
\(L'\subseteq L\), maximization over the smaller set gives
\(F_{L'}(b)\leq F_L(b)\) for every \(b\).  Coverage of
\(\mathsf{frontier}(b)\) supplies some \(s\in L'\) with
\(j_s(b)=F_L(b)\).  Therefore
\(F_L(b)=j_s(b)\leq F_{L'}(b)\), proving frontier equality pointwise.  This
also covers the zero-frontier case: \(s_0\in L'\) is itself a valid attainer.

For modules, source inclusion implies
\(U_{L'}^{\rm raw}\subseteq U_L^{\rm raw}\).  If
\(m\in U_L^{\rm raw}=C_L\), coverage of
\(\mathsf{module}(m)\) supplies a retained carrier, so
\(m\in U_{L'}^{\rm raw}\).  Thus the two raw unions are equal.  Identity
closure converts this equality into \(C_{L'}=C_L\).  A module may have one or
many source carriers; the row constraint requires at least one and does not
select a canonical carrier.
\end{proof}

\paragraph{Weighted binary formulation.}
Let the exact rational retention weights satisfy \(w_{s_0}=0\) and
\(w_s>0\) for \(s\neq s_0\).  Introduce one variable \(x_s\in\{0,1\}\) for
each \(s\in L\), where \(x_s=1\) means that \(s\) is retained.  The theorem
gives the exact formulation
\begin{align}
 \min_{x}\quad &\sum_{s\in L}w_sx_s,
     \label{eq:tagged-cover-objective}\\
 \text{s.t.}\quad
 &\sum_{s\in L:\,j_s(b)=F_L(b)}x_s\geq 1,
       && b\in B,\label{eq:tagged-frontier-rows}\\
 &\sum_{s\in L:\,m\in\mods(s)}x_s\geq 1,
       && m\in C_L,\label{eq:tagged-module-rows}\\
 &x_{s_0}=1,\qquad x_s\in\{0,1\},
       &&s\in L.\label{eq:tagged-binary-domain}
\end{align}
The source-sublibrary restriction is built into the column set \(L\).  The
selected library is \(L_x=\{s\in L:x_s=1\}\), its burden is exactly the
objective in \eqref{eq:tagged-cover-objective}, and
Theorem~\ref{thm:aor-tagged-cover-equivalence} identifies its feasible vectors
with the exact innovation-safe sublibraries.  Zero-frontier rows are
redundant after \(x_{s_0}=1\), but retaining them makes the representation a
literal tagged copy of every declared belief and avoids a hidden
positive-frontier convention.

\paragraph{Why arbitrary closure is different.}
For a nonidentity closure operator, equality of closed unions need not be
equivalent to retaining a raw carrier for each member of the source closure.
Capabilities may be generated jointly.  A correct general-closure model must
enforce the set-valued equality
\[
 \cl\!\left(\bigcup_{s:x_s=1}\mods(s)\right)=C_L
\]
or use a separately proved valid extended formulation.  The independent
module rows in \eqref{eq:tagged-module-rows} have no automatic validity in
that setting.

\paragraph{Exact boundary examples.}
The following finite examples use integer profiles and exact finite sets.

\begin{enumerate}
\item \emph{Frontier rows alone are insufficient.}  Take one belief, one
module \(m\), identity closure, and
\(L=\{s_0,s_m\}\), with both profiles zero,
\(\mods(s_0)=\varnothing\), and \(\mods(s_m)=\{m\}\).  The sublibrary
\(\{s_0\}\) covers the only frontier row and preserves the zero frontier, but
its closure is empty rather than \(\{m\}\).

\item \emph{Module rows alone are insufficient.}  Take one belief, identity
closure, and \(L=\{s_0,s_1\}\), with
\(j_{s_0}=0\), \(j_{s_1}=1\), and both module sets empty.  The source module
row family is empty, so \(\{s_0\}\) satisfies it vacuously, but its frontier
is zero rather than one.

\item \emph{Independent raw-module rows can fail under closure
complementarity.}  Let \(M=\{a,b,c\}\) and
\[
 \cl(A)=
 \begin{cases}
 A\cup\{c\},&\{a,b\}\subseteq A,\\
 A,&\text{otherwise}.
 \end{cases}
\]
This map is extensive, monotone, and idempotent.  Let the zero-profile source
be \(L=\{s_0,s_a,s_b\}\), where \(s_a\) carries only \(a\) and \(s_b\)
carries only \(b\).  Then \(C_L=\{a,b,c\}\), and the source is of course
safe-feasible relative to itself.  Nevertheless a raw-carrier row for
\(c\in C_L\) has no covering source strategy.  The naive independent-row
model therefore rejects even the source library.
\end{enumerate}

The deterministic machine-readable realization is
\path{journal/aor/fixtures/tagged_cover_equivalence_v1.json}.  It records
exact finite computation in Julia, including all sublibraries of its declared
small-library grid; it is neither an exhaustive proof over arbitrary finite
instances nor a solver optimality certificate.  The universal core
biconditional and exact burden identity are separately checked in the Lean
source \path{TaggedSafeCompression.lean} under
\path{formal/StrategyInnovation/Optimization/}.
Those declarations formalize neither a complexity result nor an
approximation guarantee.
